% Ad Familiares 5.13 --- headnote
% ad-familiares-05-13, March or April 45 BC, Astura or Ficuleanum

\textit{Cicero to Lucius Lucceius, written from Astura or
Atticus's estate at the Ficuleanum, in March or April 45 BC
(Perseus: \textit{Asturae aut in Attici Ficuleano m.~Mart.~aut
Apr.~a.~709 (45)}). Tullia, Cicero's daughter, had died at
Tusculum in mid-February; the weeks that followed are the
deepest of Cicero's grief and the period of his retreat from
Rome --- first to Astura on the coast, then to Atticus's
suburban property at the Ficuleanum --- during which he was
writing and receiving consolation letters from the friends of
his life. Lucius Lucceius, historian and senatorial intimate
(the same Lucceius whom Cicero in 56 BC had begged in
\textit{Fam.}~5.12 to write up his consulship as a monograph),
had sent one of those consolations; this is Cicero's reply.}

\textit{The letter belongs to the same consolation cluster as
the famous exchange with Servius Sulpicius (\textit{Fam.}
4.5--4.6) and the lost \textit{Consolatio} which Cicero was
writing for himself in this same period. Its register is high
and reflective: a thank-you in the philosophical mode, in
which the consoled man holds up his consoler as the better
practitioner of the doctrine they share. \S 1 praises Lucceius
for the Stoic premise that the good life cannot hang on
externals; \S 2 confesses that the storms of the last year
have shaken loose what was once deeply lodged; \S 3, the
hinge, returns the compliment of bravery with a twist --- Cicero
declares himself braver than his teacher, because where
Lucceius still hopes for the commonwealth, Cicero, looking
round at all its limbs, can find none unbroken. The closing
sections settle into the warm cadence of friendship: \S 4
recalls their shared acts of public service (the
\textit{auctore} of \S 4 alludes to Lucceius's role behind
Cicero's consular policy), and \S 5 promises that the joining
of minds will keep them together whatever distances of body and
health may impose.}
